\documentclass[10pt,a4paper]{article}

% =========================================================
% Packages
% =========================================================
\usepackage[
    a4paper,
    top=15mm,
    bottom=15mm,
    left=18mm,
    right=18mm
]{geometry}

\usepackage{fontspec}      % Compile with XeLaTeX
\usepackage{xcolor}
\usepackage{graphicx}
\usepackage{pgfcalendar}
\usepackage{tabularx}
\usepackage{array}
\usepackage{booktabs}
\usepackage{enumitem}
\usepackage[most]{tcolorbox}
\usepackage{fancyhdr}
\usepackage{float}
\usepackage{hyperref}

% =========================================================
% General formatting
% =========================================================
\setlength{\parindent}{0pt}
\setlength{\parskip}{4pt}
\renewcommand{\arraystretch}{1.20}

\definecolor{mainblue}{RGB}{38,70,110}
\definecolor{lightblue}{RGB}{238,244,250}
\definecolor{lightgray}{RGB}{246,246,246}
\definecolor{linegray}{RGB}{190,190,190}

\hypersetup{
    colorlinks=true,
    linkcolor=mainblue,
    urlcolor=mainblue
}

% =========================================================
% Automatic reporting week
%
% Week convention:
%   Monday -- Sunday
%
% Month assignment:
%   A week belongs to the month containing its Thursday.
%   This is equivalent to assigning the week to the month
%   containing at least four of its seven days.
%
% Example:
%   Mon Aug 31 -- Sun Sep 6
%   Thursday = Sep 3
%   => September, Week 1
% =========================================================

\usepackage{pgfcalendar}

% ---------------------------------------------------------
% Integer registers
% ---------------------------------------------------------
% \WeekMondayJulian:
%   Julian day corresponding to Monday of the current week.
%
% \WeekdayIndex:
%   Monday=0, ..., Sunday=6.
%
% \ScratchJulian:
%   Temporary register used for Thursday and Sunday.
%
% \ReportWeekNumber:
%   Week number within the reporting month.
% ---------------------------------------------------------

\newcount\WeekMondayJulian
\newcount\WeekdayIndex
\newcount\ScratchJulian
\newcount\ReportWeekNumber


% =========================================================
% STEP 1
% Convert compilation date into a Julian day.
%
% \year, \month, and \day are built-in TeX values.
% They correspond to the same date displayed by \today.
% =========================================================

\pgfcalendardatetojulian
    {\the\year-\the\month-\the\day}
    {\WeekMondayJulian}


% =========================================================
% STEP 2
% Determine the weekday.
%
% PGF convention:
%   Monday    = 0
%   Tuesday   = 1
%   Wednesday = 2
%   Thursday  = 3
%   Friday    = 4
%   Saturday  = 5
%   Sunday    = 6
% =========================================================

\pgfcalendarjuliantoweekday
    {\WeekMondayJulian}
    {\WeekdayIndex}


% =========================================================
% STEP 3
% Move backward to Monday.
%
% Example:
%   Saturday -> weekday index = 5
%   today - 5 days = Monday
% =========================================================

\advance\WeekMondayJulian by -\WeekdayIndex


% =========================================================
% STEP 4
% Convert Monday back to a calendar date.
% =========================================================

\pgfcalendarjuliantodate
    {\WeekMondayJulian}
    {\WeekStartYear}
    {\WeekStartMonth}
    {\WeekStartDay}


% =========================================================
% STEP 5
% Thursday = Monday + 3 days.
%
% The Thursday determines the reporting month.
% =========================================================

\ScratchJulian=\WeekMondayJulian
\advance\ScratchJulian by 3

\pgfcalendarjuliantodate
    {\ScratchJulian}
    {\ReportYear}
    {\ReportMonth}
    {\ReportThursdayDay}


% =========================================================
% STEP 6
% Determine the monthly week number.
%
% Thursday:
%
%    1--7  -> Week 1
%    8--14 -> Week 2
%   15--21 -> Week 3
%   22--28 -> Week 4
%   29--31 -> Week 5
%
% Formula:
%
%   floor((Thursday day - 1) / 7) + 1
% =========================================================

\ReportWeekNumber=\ReportThursdayDay
\advance\ReportWeekNumber by -1
\divide\ReportWeekNumber by 7
\advance\ReportWeekNumber by 1


% =========================================================
% STEP 7
% Sunday = Monday + 6 days.
%
% Reuse \ScratchJulian; Thursday is no longer needed.
% =========================================================

\ScratchJulian=\WeekMondayJulian
\advance\ScratchJulian by 6

\pgfcalendarjuliantodate
    {\ScratchJulian}
    {\WeekEndYear}
    {\WeekEndMonth}
    {\WeekEndDay}


% =========================================================
% Formatting helpers
% =========================================================

% Add a leading zero to a one-digit number.
% 8 -> 08
% 12 -> 12
\newcommand{\TwoDigits}[1]{%
    \ifnum\numexpr#1\relax<10 0\fi
    \number\numexpr#1\relax
}


% Example:
%   August 2026, Week 3
\newcommand{\ReportWeekLabel}{%
    \pgfcalendarmonthname{\ReportMonth}%
    \ \ReportYear, Week~\the\ReportWeekNumber%
}


% Example:
%   2026-08-17 -- 2026-08-23
\newcommand{\ReportWeekRange}{%
    \WeekStartYear-%
    \TwoDigits{\WeekStartMonth}-%
    \TwoDigits{\WeekStartDay}%
    \,--\,%
    \WeekEndYear-%
    \TwoDigits{\WeekEndMonth}-%
    \TwoDigits{\WeekEndDay}%
}

% ISO-formatted compilation date
\newcommand{\ReportDateISO}{%
    \CurrentYear-\CurrentMonth-\CurrentDay%
}


% =========================================================
% Header / Footer
% =========================================================
\pagestyle{fancy}
\fancyhf{}
\renewcommand{\headrulewidth}{0pt}

\fancyfoot[L]{\footnotesize Weekly Research Report}
\fancyfoot[R]{\footnotesize \thepage~/~2}


% =========================================================
% PPP box
% =========================================================
\newtcolorbox{pppbox}[1]{
    enhanced,
    breakable,
    colback=white,
    colframe=mainblue,
    boxrule=0.8pt,
    arc=1.5mm,
    left=3mm,
    right=3mm,
    top=2mm,
    bottom=2mm,
    title={#1},
    fonttitle=\bfseries\large,
    coltitle=white,
    colbacktitle=mainblue
}


% =========================================================
% Document
% =========================================================
\begin{document}

% =========================================================
% PAGE 1
% =========================================================

\begin{center}
    {\LARGE\bfseries Weekly Research Report}\\[1mm]
    {\small Progress / Problems / Plans}
\end{center}

\vspace{2mm}


% =========================================================
% Basic information
% =========================================================
\begin{tcolorbox}[
    colback=lightgray,
    colframe=linegray,
    boxrule=0.5pt,
    arc=1mm,
    left=2mm,
    right=2mm,
    top=2mm,
    bottom=2mm
]

\begin{tabularx}{\textwidth}{
    >{\bfseries}p{25mm} X
    >{\bfseries}p{30mm} X
}

Name
    & Your Name
    & Reporting Week
    & \ReportWeekLabel
    \\

Project
    & Project / Research Topic
    & Date
    & \today
    \\

Advisor
    & Advisor Name
    & Period
    & \ReportWeekRange
    \\

Main Goal
    & \multicolumn{3}{X}{
        Write the primary research objective for this week in one sentence.
      }

\end{tabularx}

\end{tcolorbox}


\vspace{3mm}


% =========================================================
% Weekly Highlight
% =========================================================
\begin{tcolorbox}[
    colback=lightblue,
    colframe=mainblue,
    boxrule=0.6pt,
    arc=1mm,
    title={Weekly Highlight},
    fonttitle=\bfseries,
    colbacktitle=mainblue
]

Summarize the most important result or conclusion of the week
in two to four sentences.

For example, the proposed loss formulation improved validation
accuracy by 2.3 percentage points over the baseline. The gain was
particularly noticeable in the low-data regime. Additional ablation
experiments will be conducted next week to determine the primary
source of this improvement.

\end{tcolorbox}


\vspace{3mm}


% =========================================================
% P1. Progress
% =========================================================
\begin{pppbox}{P1. Progress}

\textbf{1. Research Activities}

Briefly describe the main research activities performed this week.

\begin{itemize}[
    leftmargin=5mm,
    itemsep=1mm,
    topsep=1mm
]
    \item Experiment A: purpose and experimental setup.
    \item Experiment B: modified conditions or new approach.
    \item Implementation: major code or pipeline changes.
\end{itemize}


\vspace{2mm}

\textbf{2. Quantitative Results}

Summarize the primary quantitative results in the table below.

% =========================================================
% TABLE TEMPLATE
% Replace the column names and values directly.
% =========================================================
\begin{table}[H]
    \centering
    \small

    \begin{tabularx}{0.95\linewidth}{
        l
        >{\centering\arraybackslash}X
        >{\centering\arraybackslash}X
        >{\centering\arraybackslash}X
    }
        \toprule
        Method
            & Metric 1
            & Metric 2
            & Metric 3
            \\

        \midrule

        Baseline
            & 00.0
            & 00.0
            & 00.0
            \\

        Proposed A
            & 00.0
            & 00.0
            & 00.0
            \\

        Proposed B
            & \textbf{00.0}
            & \textbf{00.0}
            & \textbf{00.0}
            \\

        \bottomrule
    \end{tabularx}

    \caption{
        Main experimental results.
        Replace this text with a concise description of the table.
    }
    \label{tab:main-results}
\end{table}


Write two or three sentences interpreting the table.
Focus on the main observation rather than repeating every number.

\end{pppbox}


\vspace{2mm}


% =========================================================
% FIGURE TEMPLATE
%
% Put the actual figure at:
%   figures/main-result.pdf
%
% PDF is recommended for plots/vector graphics.
% PNG/JPG also work.
%
% If the file does not exist, a placeholder is shown,
% so this template still compiles.
% =========================================================
\begin{figure}[H]
    \centering

    \IfFileExists{figures/main-result.pdf}
    {
        \includegraphics[
            width=0.72\linewidth
        ]{figures/main-result.pdf}
    }
    {
        % Placeholder shown before you add the actual figure.
        \fbox{
            \parbox[c][35mm][c]{0.72\linewidth}{
                \centering
                \textbf{Figure Placeholder}\\[2mm]
                Replace with\\
                \texttt{figures/main-result.pdf}
            }
        }
    }

    \caption{
        Main experimental result.
        State what is plotted and the key experimental condition.
    }

    \label{fig:main-result}
\end{figure}


% =========================================================
% PAGE 2
% =========================================================
\newpage


% =========================================================
% P2. Problems
% =========================================================
\begin{pppbox}{P2. Problems}

\textbf{1. Main Issue}

Describe the most important unresolved issue or bottleneck.

For example, training remains unstable on a particular dataset,
with substantial variation across random seeds. It is currently
unclear whether this behavior originates from the learning rate,
data distribution, or optimization procedure.


\vspace{2mm}

\textbf{2. Working Hypotheses}

\begin{itemize}[
    leftmargin=5mm,
    itemsep=1mm,
    topsep=1mm
]
    \item Hypothesis 1: possible explanation and supporting evidence.
    \item Hypothesis 2: alternative explanation.
    \item Hypothesis 3: additional factor that should be tested.
\end{itemize}


\vspace{2mm}

\textbf{3. Attempts Made}

List approaches already attempted and their outcomes.

\begin{itemize}[
    leftmargin=5mm,
    itemsep=1mm,
    topsep=1mm
]
    \item Attempt 1 $\rightarrow$ observed outcome.
    \item Attempt 2 $\rightarrow$ observed outcome.
\end{itemize}


\vspace{2mm}

\textbf{4. Questions for Discussion}

\begin{itemize}[
    leftmargin=5mm,
    itemsep=1mm,
    topsep=1mm
]
    \item Question or decision to discuss with the advisor.
    \item Additional question requiring feedback.
\end{itemize}

\end{pppbox}


\vspace{3mm}


% =========================================================
% P3. Plans
% =========================================================
\begin{pppbox}{P3. Plans}

\textbf{1. Objectives for Next Week}

\begin{enumerate}[
    leftmargin=6mm,
    itemsep=1mm,
    topsep=1mm
]
    \item \textbf{High:}
          Perform the highest-priority experiment or analysis.

    \item \textbf{High:}
          Conduct an additional experiment addressing the main issue.

    \item \textbf{Medium:}
          Implement supporting code or review relevant literature.

    \item \textbf{Low:}
          Perform optional analysis if time permits.
\end{enumerate}


\vspace{2mm}

\textbf{2. Execution Plan}

% =========================================================
% PLANNING TABLE TEMPLATE
% =========================================================
\begin{table}[H]
    \centering
    \small

    \begin{tabularx}{0.96\linewidth}{
        >{\centering\arraybackslash}p{18mm}
        X
        >{\centering\arraybackslash}p{31mm}
    }
        \toprule
        Priority
            & Task
            & Expected Output
            \\

        \midrule

        High
            & Run Experiment A with three random seeds.
            & Result table
            \\

        High
            & Perform an ablation study on component B.
            & Ablation plot
            \\

        Medium
            & Review two or three related papers.
            & Literature notes
            \\

        \bottomrule
    \end{tabularx}

    \caption{Planned tasks for the next reporting period.}
    \label{tab:next-week}
\end{table}


\textbf{3. Expected Outcome}

State the concrete result that should be available before the next
weekly meeting.

For example, complete experiments with at least three random seeds
and prepare a comparison table and plot showing the difference
between the baseline and the proposed method.

\end{pppbox}


\vfill


% =========================================================
% Meeting notes
% =========================================================
\begin{tcolorbox}[
    colback=lightgray,
    colframe=linegray,
    boxrule=0.5pt,
    arc=1mm,
    title={Meeting Notes / Feedback},
    fonttitle=\bfseries,
    coltitle=black,
    colbacktitle=lightgray
]

\vspace{14mm}

% Write advisor / lab feedback here during the meeting.

\end{tcolorbox}


\end{document}